\documentclass[12pt]{article}

\usepackage{amsmath}
\usepackage{amssymb}
\usepackage{graphicx}
\usepackage{hyperref}
\usepackage[margin=1in]{geometry}
\usepackage{natbib}
\usepackage{xcolor}
\usepackage{tcolorbox}
\usepackage{booktabs}
\usepackage{tabularx}
\usepackage{longtable}
\usepackage{lineno}
\linenumbers
\tcbuselibrary{breakable,skins}
\bibliographystyle{aasjournal}

\newcommand{\cmg}{\,\mathrm{cm^2\,g^{-1}}}
\newcommand{\kms}{\,\mathrm{km\,s^{-1}}}
\newcommand{\Msun}{\,M_\odot}
\newcommand{\smx}{\sigma/m_\chi}
\newcommand{\sigmam}{\sigma/m}
\newcommand{\vrel}{v_{\rm rel}}
\newcommand{\Kn}{\mathrm{Kn}}
\newcommand{\mfp}{\lambda_{\rm mfp}}

\title{Self-Interacting Dark Matter}
\author{}
\date{\today}

\begin{document}

\maketitle
\tableofcontents
\newpage

\begin{tcolorbox}[colback=yellow!10!white, colframe=orange!75!black,
                  title={\bfseries Disclaimer}]
These notes were compiled by an AI assistant (Claude, Anthropic) and have
\textbf{not been reviewed by a human domain expert}.  They may contain
errors in derivations, physical reasoning, or numerical estimates.
Cross-check key equations against the primary literature before use.
\end{tcolorbox}
\vspace{0.5em}

%----------------------------------------------------------------------
\section{Motivation}
\label{sec:motivation}

Cold dark matter (CDM) treats DM as a collisionless species, and this
description reproduces cosmological structure at the few-percent level
on scales above $\sim\!1\,{\rm Mpc}$.  On sub-galactic scales three
tensions historically motivated the SIDM proposal, all mediated by the
central density structure of DM halos:

\begin{itemize}
  \item \textbf{Core--cusp.}  Rotation curves and stellar velocity
        dispersions in low-surface-brightness galaxies and dwarf
        spheroidals prefer approximately constant-density inner
        profiles, while CDM $N$-body simulations produce $\rho\propto
        r^{-1}$ NFW cusps.
  \item \textbf{Too-big-to-fail.}  The most massive Milky Way
        subhaloes predicted by CDM are too concentrated at their
        centres to host the observed classical dwarfs.
  \item \textbf{Diversity of rotation curves.}  Nearby dwarfs of
        similar circular velocity show a wider spread of inner slopes
        than CDM sims produce.
\end{itemize}

\citet{SpergelSteinhardt2000} proposed resolving these tensions by
giving dark matter an elastic self-interaction with cross section per
unit mass $\smx$ of order $1\cmg$: large enough to matter in dense
environments (galactic cores) and negligible on cosmological scales.
The corresponding mean free path is $1\,{\rm kpc}\lesssim
\mfp\lesssim 1\,{\rm Mpc}$ at the solar-neighbourhood density
$\rho\simeq 0.4\,{\rm GeV\,cm^{-3}}$, requiring
\citep{SpergelSteinhardt2000}
%
\begin{equation}
  0.45\cmg \le \smx \le 450\cmg,
  \qquad
  1\,{\rm MeV} \le m_\chi \le 10\,{\rm GeV}
\end{equation}
%
in the original paper's ``elastic, non-dissipative, non-annihilating''
setup.  Twenty-five years of $N$-body simulation, dwarf-galaxy
kinematics, and cluster observations have collapsed this window from
three decades to a factor of $\sim\!10$, and have made clear that a
single constant $\smx$ across all scales is disfavoured.  What
survives is a \emph{velocity-dependent} $\sigma(v)$, running from
$\sim\!100\cmg$ on dwarf scales to $\sim\!0.1\cmg$ on cluster scales
--- naturally realised by a light-mediator Yukawa dark force
\citep{Kaplinghat2016,TulinYu2018}.

The modern picture (\S\ref{sec:regimes}--\ref{sec:constraints}) has
also shifted: the phenomenologically important process is not just
\emph{core formation} (heat conduction outward flattens the cusp) but
also \emph{core collapse} (gravothermal catastrophe drives the centre
into a runaway compact state).  Both branches are now recognised as
essential to explaining the observed diversity of subhaloes
\citep{Zeng2022,Nadler2023}, and possibly to seeding the observed
$z\gtrsim 7$ supermassive black holes
\citep{Pollack2015,Feng2021}.

%----------------------------------------------------------------------
\section{Physics of the self-interaction}
\label{sec:physics}

\subsection{Cross section, mean free path, scattering rate}

For an elastic cross section $\sigma$ and DM particle mass $m_\chi$,
the local mean free path in a halo of density $\rho$ is
%
\begin{equation}
  \mfp = \frac{m_\chi}{\rho\,\sigma} = \frac{1}{\rho\,\smx},
\end{equation}
%
and the corresponding scattering rate per particle is
%
\begin{equation}
  R_{\rm scat} = \rho\,\smx\,v_{\rm rel}
    \approx 0.1\,{\rm Gyr}^{-1}
      \left(\frac{\rho}{0.1\,M_\odot\,{\rm pc}^{-3}}\right)
      \left(\frac{\vrel}{50\kms}\right)
      \left(\frac{\smx}{1\cmg}\right).
  \label{eq:Rscat}
\end{equation}
%
Normalisation \eqref{eq:Rscat} makes the astrophysical target
explicit: for typical dwarf-halo densities and virial velocities,
$\smx \sim 1\cmg$ delivers order-one scatterings per particle per
Hubble time --- enough to redistribute energy in the core, negligible
for the outer halo.  Standard WIMP self-scattering rates are
$\smx \sim 10^{-38}\,{\rm cm^2\,GeV^{-1}}$, some $14$ orders of magnitude
below the astrophysical target \citep{TulinYu2018}, so SIDM is
manifestly \emph{not} a WIMP.

Units conversion:
$1\cmg = 1.78\times 10^{-24}\,{\rm cm^2\,GeV^{-1}}$.  In particle-physics
Born-limit terms, a Yukawa mediator of mass $m_\phi$ coupled to DM
with strength $\alpha'$ gives
%
\begin{equation}
  \smx \simeq \frac{4\pi\alpha'^2 m_\chi}{m_\phi^4}
    \approx 1\cmg
      \left(\frac{\alpha'}{0.01}\right)^{\!2}
      \left(\frac{m_\chi}{10\,{\rm GeV}}\right)
      \left(\frac{m_\phi}{40\,{\rm MeV}}\right)^{\!-4}.
  \label{eq:Yukawa_Born}
\end{equation}
%
Electromagnetic-strength dark coupling ($\alpha' \sim 10^{-2}$) with a
sub-GeV mediator produces the astrophysical target; the DM particle
itself does not need to be strongly coupled.  That last statement
separates two meanings of ``strong'' that are easy to conflate:
strongly \emph{interacting} in the astrophysical sense (a large cross
section per unit mass, $\smx\sim 1\cmg$) versus strongly
\emph{coupled} in the particle-physics sense (a large dimensionless
coupling $\alpha'\gtrsim 1$, non-perturbative dynamics).  SIDM needs
the first, not the second.  Equation~\eqref{eq:Yukawa_Born} shows
why: the cross section scales as the inverse \emph{fourth} power of
the mediator mass, so the enormous $\sigma$ comes from the mediator
being light ($m_\phi\sim 10$--$40\,{\rm MeV}$, force range of tens of
fm --- long-range by particle standards), not from the coupling being
large.  Taking the same $\alpha' = 10^{-2}$ but a weak-scale mediator
$m_\phi = 100\,{\rm GeV}$ suppresses $\sigma$ by
$(100\,{\rm GeV}/40\,{\rm MeV})^4 \approx 4\times 10^{13}$ ---
essentially the entire 14 orders of magnitude separating SIDM from
WIMP self-scattering is mediator mass, none of it coupling strength.
Familiar analogues: Rutherford scattering, whose huge low-energy
cross sections come from a massless photon at $\alpha = 1/137$; and
nucleon--nucleon scattering, where $\sigma/m_n \sim$ a few $\cmg$
--- comparable to the SIDM target --- is set by the pion's range, not
by any coupling being infinite.  The model-building consequence:
$\alpha'$ stays perturbative, so the Born calculation is valid, a
dark photon plus dark Higgs suffices as the underlying theory, and
the Rutherford-like $\sigma(v)\propto v^{-4}$ falloff at high
velocity --- the velocity dependence the data prefer
(\S\ref{sec:vdep}) --- emerges automatically rather than being
engineered.

\subsection{Two regimes: cored and core-collapsed}
\label{sec:regimes}

Whether a halo is in the ``dilute'' or ``dense'' regime depends on how
the accumulated scattering count $R_{\rm scat}\, t_{\rm age}$ compares
with unity in its centre.  A convenient dimensionless proxy is the
Knudsen number
%
\begin{equation}
  \Kn \equiv \frac{\mfp}{H},
  \qquad
  H \equiv \sqrt{\frac{v^2}{4\pi G\rho}},
\end{equation}
%
where $H$ is the gravitational scale height.  $\Kn\gg 1$ is the
long-mean-free-path (LMFP) regime, appropriate for the outer halo and
the cored phase.  $\Kn\ll 1$ is the short-mean-free-path (SMFP)
regime, appropriate for the collapsing centre.

\paragraph{Core-forming (LMFP) regime.}  When scatterings deliver
\emph{outward} heat conduction from the denser centre, the inner
density profile is flattened into an approximately isothermal core.
The core radius $r_c$ is where $\mfp\sim r_c$; typical values are
$r_c \sim 1$--$3\,{\rm kpc}$ for dwarf-scale $\smx \sim 1\cmg$
\citep{Rocha2013}.  A stable equilibrium is reached in
$\sim 1$--$10\,{\rm Gyr}$, well before core collapse.

\paragraph{Core-collapsed (SMFP) regime.}  When the accumulated
scattering count in the centre exceeds unity, the core becomes a
tightly self-coupled fluid.  \emph{Gravothermal catastrophe} sets in:
the negative heat capacity of a self-gravitating system means outward
conduction of heat cools the centre --- but the centre then contracts,
heating up further, so the process runs away
\citep{Balberg2002}.  The central density rises without bound, the
inner profile transitions from a flat core through the LMFP
self-similar $\rho\propto r^{-2.19}$ solution and into a smaller,
denser SMFP inner core.  Fluid-model calculations give a universal
collapse time
%
\begin{equation}
  t_{\rm coll} \simeq 290\,t_r
  \qquad\text{(LMFP self-similar, hard-sphere)}
  \qquad\text{\citep{Balberg2002}}
  \label{eq:tcoll_290}
\end{equation}
%
or with more modern calibration
$t_{\rm coll} \simeq 382\,t_r$ \citep{Meshveliani2023}
or
$t_{\rm coll} \simeq 456\,t_r$ \citep{Pollack2015},
where the SIDM relaxation time is
%
\begin{equation}
  t_r = \frac{1}{a\,\rho\,\sigma\,v},
  \qquad
  a = \sqrt{16/\pi}\simeq 2.26,
  \label{eq:tr_def}
\end{equation}
%
for hard-sphere scattering with Maxwellian velocities.  For typical
dwarf conditions ($\rho\sim 10^{-24}\,{\rm g\,cm^{-3}}$, $v\sim
10^7\,{\rm cm\,s^{-1}}$, $\sigma\sim 1\cmg$), $t_r\sim 1.4\,{\rm Gyr}$
so $t_{\rm coll}\gtrsim 10\,t_{\rm Hubble}$ --- present-day dwarf cores
are safe from collapse at $\smx\sim 1\cmg$
\citep{Balberg2002}.  Collapse can be forced within a Hubble time only
by cranking $\smx$ into the tens or hundreds of $\cmg$
\citep{Correa2021,Meshveliani2023}, or by shortening $t_r$ via baryon
compression (\S\ref{sec:baryons}) or tidal stripping.

\citet{Zeng2022} calibrated the LMFP fluid prediction against
$N$-body sims and reported the empirical scaling
%
\begin{equation}
  t_{\rm coll} \simeq 3.58\times 10^{8}\,{\rm Gyr}
    \left(\frac{\smx}{\cmg}\right)^{-0.74}
    \left(\frac{M_{200m}}{M_\odot}\right)^{-0.24}
    c_{200m}^{-2.67},
  \label{eq:zeng_tcoll}
\end{equation}
%
shallower in $\smx$ and $c$ than the fluid prediction
$t_c\propto\smx^{-1}c^{-7/2}$, reflecting numerical effects that the
fluid model does not capture.

\subsection{Velocity dependence and the Yukawa dark force}
\label{sec:vdep}

Data across scales require a velocity-dependent $\sigma(v)$
\citep{Kaplinghat2016,TulinYu2018}.  In the Born limit, a Yukawa
mediator gives a Rutherford-like differential cross section commonly
parameterised as
%
\begin{equation}
  \frac{d\sigma}{d\cos\theta}
    = \frac{\sigma_0\,w^4}{2\bigl[w^2 + v^2\sin^2(\theta/2)\bigr]^2},
  \label{eq:rutherford_sidm}
\end{equation}
%
with $\sigma_0$ the low-velocity amplitude and $w$ the transition
velocity above which the cross section drops as $v^{-4}$.
For $\alpha'\sim 0.01$, $m_\chi\sim 10\,{\rm GeV}$, $m_\phi\sim
40\,{\rm MeV}$ one recovers $\smx\sim 1\cmg$ on galactic scales while
falling well below cluster bounds
\citep{Kaplinghat2016,TulinYu2018}.  Modern cosmological-scale
zoom-in suites use benchmark ``Group SIDM'' or ``MW SIDM'' models with
$\sigma_0/m = 147\cmg$ and $w = 24$ or $120\kms$
\citep{Nadler2023,Palubski2025}, chosen so that $\sigma_{\rm eff}/m$
reaches $\gtrsim 100\cmg$ on dwarf scales (driving core collapse in
concentrated subhaloes) while dropping below $1\cmg$ on cluster
scales.

For pipeline calibration it is useful to define a viscosity-weighted
Maxwellian-averaged effective cross section
%
\begin{equation}
  \sigma_{\rm eff}
    = \frac{2\!\int v^2\,dv\,d\cos\theta\;
      \frac{d\sigma}{d\cos\theta}\sin^2\theta\,v^5\,
      e^{-v^2/(4\nu_{\rm eff}^2)}}
      {\int v^2\,dv\,d\cos\theta\;
      \sin^2\theta\,v^5\,e^{-v^2/(4\nu_{\rm eff}^2)}},
  \qquad
  \nu_{\rm eff}\simeq 0.64\,V_{\rm max},
  \label{eq:sigma_eff}
\end{equation}
%
which collapses the angular and velocity distributions into one number
per halo \citep{Nadler2023}.

%----------------------------------------------------------------------
\section{Gravothermal fluid framework}
\label{sec:fluid}

The workhorse semi-analytic model for isolated SIDM halos is the
gravothermal conducting-fluid formalism of
\citet{Balberg2002}, adapted from stellar-dynamical globular-cluster
theory.  For readers unfamiliar with the kinetic-theory chain
Boltzmann $\to$ moments $\to$ Chapman--Enskog closure $\to$
Navier--Stokes, Appendix~\ref{app:kinetic} gives the derivation from
scratch.  What follows uses those results without proof.

For a spherical halo with 1D velocity dispersion $\nu$ and
luminosity $L$, one solves
%
\begin{align}
  \frac{\partial M}{\partial r} &= 4\pi r^2\rho,
  \label{eq:balberg_mass}\\
  \frac{\partial(\rho\nu^2)}{\partial r} &= -\frac{GM\rho}{r^2},
  \label{eq:balberg_hse}\\
  \frac{\partial L}{\partial r} &= -4\pi r^2\rho\nu^2
    \left(\frac{\partial}{\partial t}\right)_{\!M}\!\ln\frac{\nu^3}{\rho},
  \label{eq:balberg_energy}\\
  \frac{L}{4\pi r^2}
   &= -\frac{3}{2}ab\,\nu\sigma
      \left[a\sigma^2 + \frac{4\pi G b}{C\rho\nu^2}\right]^{-1}
      \frac{\partial\nu^2}{\partial r},
  \label{eq:balberg_flux}
\end{align}
%
where $a = \sqrt{16/\pi}$, $b = 25\sqrt{\pi}/32 \simeq 1.0$ from
Chapman--Enskog, and $C$ is an order-unity interpolation constant
calibrated against $N$-body: $C \simeq 0.75$
\citep{Balberg2002,Pollack2015} or $C \simeq 0.42$ from a recent
velocity-dependent-SIDM $N$-body match \citep{Meshveliani2023}.
Equation~\eqref{eq:balberg_flux} interpolates between the SMFP and LMFP
conductivities:
%
\begin{equation}
  \kappa_{\rm SMFP} \simeq 2.1\,\nu k_B/\sigma,
  \qquad
  \kappa_{\rm LMFP} \simeq 0.27\,C\rho\nu^3\sigma k_B/(Gm^2),
\end{equation}
%
and reduces to Fickian conduction $L \propto \kappa\,dT/dr$ in each
limit.  Integrating
\eqref{eq:balberg_mass}--\eqref{eq:balberg_flux} on a Lagrangian mass
grid produces the characteristic ``core forms then collapses'' history
of an isolated SIDM halo, ending in a singular state.

\paragraph{Limits of the fluid model.}  The fluid model is an
adiabatic transport approximation --- it assumes local thermodynamic
equilibrium and isotropic velocity dispersion.  \citet{Gurian2025Core}
performed the first direct simulation Monte Carlo (DSMC) calculation
of gravothermal collapse (code KiSS-SIDM, publicly at
\url{https://gitlab.com/Socob/KiSS-SIDM}) and showed that
the intermediate-MFP region just outside the collapsing core stays
persistently out of LTE: the radial velocity dispersion develops
non-Maxwellian tails from ``boil-off'' particles, and the fluid
model's temperature-weighted entropy production $\hat Q$ turns positive
--- an unambiguous signal of failure.  The local
$d\log M_{\rm core}/d\log\langle v^2\rangle$ exponent in the SMFP
regime is $-0.21$ (DSMC) versus $-0.27$ to $-0.37$ (fluid) --- an
order-one disagreement in an exponential scaling that compounds over
the many $e$-folds remaining to the relativistic instability --- so
the seed mass at relativistic instability could be substantially
larger than fluid extrapolations predict.  Most of the halo's lifetime is still
LMFP-dominated, so collapse-time corrections are $\sim 30\%$ --- but
current SIDM $N$-body codes already disagree at that level, and DSMC is
now the reference calibration point \citep{Gurian2025Core}.

%----------------------------------------------------------------------
\section{Halo phenomenology}
\label{sec:halo}

\subsection{Isothermal core and the isothermal-Jeans model}

Inside the cored phase, the SIDM equilibrium is well described by an
isothermal Jeans model:
%
\begin{equation}
  \sigma_0^2\,\nabla^2\ln\rho = -4\pi G(\rho+\rho_b),
  \label{eq:isothermal_jeans}
\end{equation}
%
with $\sigma_0$ the constant 1D velocity dispersion, $\rho_b$ the
baryon density, and the outer profile matched onto an NFW envelope at
the radius $r_1$ where the accumulated scattering count equals unity:
%
\begin{equation}
  \frac{\langle\sigma v\rangle}{m}\,\rho(r_1)\,t_{\rm age} \simeq 1.
  \label{eq:r1_def}
\end{equation}
%
\citet{Kaplinghat2016} fit this two-parameter model
(inner $\rho_0$, $\sigma_0$; NFW parameters $\rho_s, r_s$ outside) to
12 dwarf+LSB rotation curves and 6 relaxed clusters, spanning
$\vrel$ from $\sim 30\kms$ to $\sim 1500\kms$.  The result:
%
\begin{align}
  \smx &\simeq 1.9^{+0.6}_{-0.4}\cmg
        & \text{(galaxies, }\vrel \sim 30\text{--}200\kms\text{)},
   \\
  \smx &= 0.10^{+0.03}_{-0.02}\cmg
        & \text{(clusters, }\vrel \sim 1500\kms\text{)}.
\end{align}
%
The two values differ by more than an order of magnitude, ruling out
constant $\smx$.  Fit to the Yukawa cross section
\eqref{eq:rutherford_sidm} (with $\alpha' = \alpha_{\rm EM}$) gives
$m_\chi = 15^{+5}_{-3}\,{\rm GeV}$ and $m_\phi = 17\pm 4\,{\rm MeV}$
\citep{Kaplinghat2016}.

\subsection{Baryonic backreaction}
\label{sec:baryons}

For DM-only ($\sigma$-only) simulations, \citet{Rocha2013} found the
SIDM core radius scales as a fixed fraction of the underlying NFW
scale radius:
%
\begin{equation}
  r_b/r_s \simeq 0.71\quad\text{for }\smx = 1\cmg,
  \qquad
  r_b/r_s \simeq 0.20\quad\text{for }\smx = 0.1\cmg,
\end{equation}
%
across four decades in halo mass.  Values of $\smx = 1\cmg$ overshoot
cluster core densities by an order of magnitude, while
$\smx = 0.1\cmg$ (velocity-independent!) matches from dwarfs to
clusters --- an existence proof for a workable constant-$\sigma$ pipeline
\citep{Rocha2013}.

Baryons dramatically alter this picture in baryon-dominated galaxies.
\citet{Kaplinghat2014} solved the isothermal-Jeans equation
\eqref{eq:isothermal_jeans} for the Milky Way (Hernquist bulge plus
disc) and found the SIDM core radius shrinks to
%
\begin{equation}
  r_c \simeq 0.3\,{\rm kpc}
    \left(\frac{r_0}{2.7\,{\rm kpc}}\right)
    \left(\frac{\sigma_0}{150\kms}\right)^{\!2}
    \left(\frac{183\kms}{V_B(r_0)}\right)^{\!2},
  \label{eq:MW_core}
\end{equation}
%
with central density $\rho_0 \sim 14$--$80\,{\rm GeV\,cm^{-3}}$ (with
and without adiabatic contraction).  Both are more than an order of
magnitude denser/smaller than the DM-only prediction --- baryons pull
the SIDM core inward, and the SIDM density tracks the deeper
baryonic potential rather than sitting in a large, dilute core.
Consequences: (i) SIDM indirect-detection signals from the Galactic
Centre are \emph{not} suppressed; (ii) cluster/Bullet bounds should be
loosened; (iii) the SIDM density inherits the non-spherical baryonic
isopotential shape.

Independent confirmation from full FIRE-2 hydrodynamic zoom-ins of
four $M_{\rm halo}\sim 10^{10}\Msun$ dwarfs \citep{Robles2017}: for
$\smx=1\cmg$, SIDM density profiles are \emph{robust to baryonic
feedback} at $M_\star < 3\times 10^6\Msun$ (SIDM-DMO and SIDM-Hydro
essentially agree), whereas CDM cores form only for
$M_\star > 3\times 10^6\Msun$ via strong stellar feedback.
Consequence: ultra-faint dwarfs with $M_\star \lesssim 3\times
10^6\Msun$ are the cleanest observational discriminator ---
feedback-free enough that CDM stays cuspy and SIDM predictions are
firm.

\subsection{Subhalo core collapse and the diversity problem}
\label{sec:subhalos}

The Milky Way subhalo population appears to require \emph{both}
core-formed and core-collapsed objects.  \citet{Correa2021} fits nine
bright MW dSphs individually and finds each of them must currently be
in a gravothermal collapse phase, with the required $\smx$ ranging
%
\begin{equation}
  \smx \simeq 20\text{--}30\cmg \text{ (Draco)}
    \quad \text{ to } \quad
  \smx \simeq 90\text{--}150\cmg \text{ (LeoII, Sextans)},
\end{equation}
%
increasing at lower orbital pericentre (lower in-core $\vrel$).  A
single Yukawa model ($m_\chi = 0.648\,{\rm GeV}$, $m_\phi =
0.636\,{\rm MeV}$, $\alpha_\chi = 0.01$) simultaneously fits all nine
plus extrapolates below cluster bounds.  Tidal stripping accelerates
gravothermal collapse by factor $\sim 10$; without it, none of the
dSphs would have collapsed by $z=0$.

\citet{Zeng2022} tempered this picture with detailed subhalo sims and
found that constant $\smx$ makes cosmological subhaloes essentially
unable to core-collapse: the critical initial concentration for a
$10^{10.5}\Msun$ subhalo to collapse in 10 Gyr is
$c\gtrsim 40\text{--}60$, well above the cosmological median $c\sim
24$.  This is because host-subhalo evaporation and tidal heating fight
the internal collapse.  Constant-$\sigma$ SIDM is therefore
disfavoured for producing the observed compact-substructure population
--- velocity dependence is required.

\citet{Nadler2023} settled this with the first cosmological zoom-in
of \emph{strong} velocity-dependent SIDM on a group-scale
($10^{13}\Msun$) host, using the ``Group SIDM'' benchmark
$\sigma_0/m = 147.1\cmg$, $w = 120\kms$.  Both extremes of the
observed subhalo population emerge from the \emph{same} model: (i) a
core-collapsed subhalo matching the dense strong-lens perturber in
SDSSJ0946+1006 (inner projected mass $2\times 10^9\Msun$ within 1 kpc,
$\gtrsim 3\sigma$ denser than CDM), and (ii) core-expanded
isolated haloes matching gas-rich UDG rotation curves.  The effective
cross section spans $\sigma_{\rm eff}/m \sim 2\text{--}100\cmg$ over
the relevant $V_{\rm max}$ range while remaining $\lesssim 1\cmg$ at
cluster scales.

The \citet{Palubski2025} SIDM Concerto public data release
extends this to 14 zoom hosts from LMC to low-mass cluster scale.  The
core-collapsed-fraction distribution peaks where the cross section
transitions from $v^{-4}$ (high velocity) to $v^0$ (low velocity):
$V_{\rm max}\sim 30\kms$ for isolated haloes, $\sim 60\kms$ for
subhaloes.  Cluster-scale hosts show no core-collapsed subhaloes,
consistent with $\smx\lesssim 0.1$--$1\cmg$ cluster bounds.

%----------------------------------------------------------------------
\section{Numerical methods}
\label{sec:numerical}

This section is a survey of the method landscape.  The companion
note \texttt{sidm\_sim.tex} develops the simulation program in
depth: the DSMC algorithm with derivations, the 3D generalization
and its cost, moment-closure schemes (a two-moment gravothermal
model), asymptotic-preserving integrators, and an implementation
plan for the TIGRIS code.  The division of labor: this note covers
the phenomenon; \texttt{sidm\_sim.tex} covers simulating it.

\subsection{Particle-based Monte Carlo}

The Rocha--Peter--Bullock--Kaplinghat scheme \citep{Rocha2013}, now
standard in GADGET, GIZMO, and Arepo, evaluates a pair-wise scattering
probability
%
\begin{equation}
  P_{ij} = \rho_{ij}\,|\mathbf{v}_i - \mathbf{v}_j|\,\smx\,\Delta t,
  \qquad
  \rho_{ij}
    = m_p\!\int d^3r\,W(|\mathbf{r}-\mathbf{r}_i|,h)
                    W(|\mathbf{r}-\mathbf{r}_j|,h),
  \label{eq:Pij}
\end{equation}
%
with $W$ a cubic-spline kernel and $h\gtrsim 0.2\times$ mean particle
spacing.  Pairs with $P_{ij}$ times a random draw over threshold scatter
elastically in the centre-of-mass frame (uniform on the sphere for
isotropic $\sigma$; drawn from
Eq.~\eqref{eq:rutherford_sidm} for Yukawa).  Momentum and energy are
conserved per pair.

The main numerical challenge is the timestep: at high central density
during collapse, $P_{ij}$ grows and the SIDM timestep must shrink
below the gravitational one, dominating cost.  Modern implementations
attempt to save on cost via adaptive timesteps and quiet mesh
refinement, but pure $N$-body sims have historically stopped at
$\rho_{\rm cen}\sim 10^3\rho_s$ \citep{Palubski2025}, well before the
relativistic instability.

\subsection{Fluid / gravothermal semi-analytic}

Fast, spherically symmetric, cheap: solve
Eqs.~\eqref{eq:balberg_mass}--\eqref{eq:balberg_flux} on a Lagrangian
mass grid.  In-house implementations accompany most of the
fluid-model literature
\citep{Balberg2002,Pollack2015,Correa2021,Meshveliani2023},
differing mainly in the calibration of the conductivity constant
$C$ against $N$-body simulations.  Limitations:
spherical only, LTE assumed, breaks at IMFP boundary
\citep{Gurian2025Core}.

\subsection{Direct simulation Monte Carlo (KiSS-SIDM)}

\citet{Gurian2025Core} present the first spherically-symmetric DSMC
kinetic solver for SIDM, using the ``no-time-counter'' scattering
scheme (Bird 1994) on an adaptive Eulerian radial grid with cells kept
below the local mean free path and Jeans length.  With $2\times 10^6$
tracer particles, KiSS-SIDM conserves total energy to
$|\Delta E/E| < 2\times 10^{-4}$ and reproduces Chapman--Enskog SMFP
conductivity within 2\%.  The full fiducial run --- an NFW halo with
$\smx\!\cdot\!\rho_s r_s = 0.32$ --- reaches central density
$\rho_c/\rho_s \sim 10^5$ in $\sim 30$ minutes on a laptop.  The main
scientific output is the identification of persistent LTE breakdown
in the IMFP region.

\subsection{Hybrid particle+fluid (SHH)}

\citet{Schon2025} introduce the SHH (SIDM Hydro Hybrid) method: couple
the standard $N$-body SIDM scattering algorithm to smoothed particle
hydrodynamics through a density-dependent bridging function
%
\begin{equation}
  h(\rho) = \left[1 + (\rho_{\rm SHH}/\rho)^{n_{\rm SHH}}\right]^{-1},
  \label{eq:SHH_h}
\end{equation}
%
with the momentum update
%
\begin{equation}
  \frac{du_i}{dt}
    = h(\rho)\,{\rm Kick}_{\rm hydro}
    + \bigl(1-h(\rho)\bigr)\,{\rm Kick}_{\rm SIDM}
    - \partial_i\Phi.
  \label{eq:SHH_momentum}
\end{equation}
%
The threshold $\rho_{\rm SHH}$ is calibrated so that hydro takes over
at $\Kn_{\rm sim}\lesssim 10^{-3}$.  In an isolated $1.15\times
10^9\Msun$ halo with $\smx = 50\cmg$, SHH reaches central density
$\rho_c\sim 10^{4-5}\rho_s$ --- two orders of magnitude beyond
pure-SIDM at fixed particle count, and in less wall time because the
runaway SIDM timestep shrinkage is replaced by cheap hydro pressure
kicks.  The implementation is ideal-fluid only at present; viscosity
and thermal conduction (Chapman--Enskog non-ideal terms) are on the
roadmap.

\subsection{Data release: SIDM Concerto}

\citet{Palubski2025} released a public library of 14 cosmological
zoom-in sims covering LMC ($10^{11}\Msun$), Milky Way
($10^{12}\Msun$), group ($10^{13}\Msun$), and low-mass cluster
($10^{14}\Msun$) hosts, in CDM and three velocity-dependent SIDM
variants (GroupSIDM, GroupSIDM-70, MilkyWaySIDM).  Each host has
$\sim 2\times 10^7$ particles;
$m_p$ and softening are 6.3$\times 10^3\Msun$, 40 pc/$h$ (LMC) to
$2.7\times 10^7\Msun$, 600 pc/$h$ (cluster).  Data DOI:
\texttt{10.5281/zenodo.14933624}.

%----------------------------------------------------------------------
\section{Observational constraints}
\label{sec:constraints}

Assembled bounds and preferences from recent literature (numbers from
\citealt{Kaplinghat2016,TulinYu2018,Nadler2023,Palubski2025}):

\begin{tabularx}{\linewidth}{@{}l l X@{}}
\toprule
Probe & $\vrel\;[\kms]$ & $\smx\;[\cmg]$ \\
\midrule
Bullet cluster (mass/light offset)          & 4000
   & $\lesssim 0.7$ \\
Bullet cluster (subcluster survival)        & 2000
   & $\lesssim 1.25$ (68\% CL) \\
Ensemble merging clusters                   & 900
   & $\lesssim 0.47$ (95\% CL) \\
Cluster halo shape / ellipticity            & 1300
   & $\lesssim 1$ \citep{TulinYu2018} \\
Cluster core density (relaxed)              & 1500
   & $= 0.10^{+0.03}_{-0.02}$ \citep{Kaplinghat2016} \\
Group core (Nadler+2023 host)               & 800
   & $\sigma_{\rm eff}/m \lesssim 1$ \\
Galaxy rotation curves (dwarfs, LSBs)       & 30--200
   & $= 1.9^{+0.6}_{-0.4}$ \citep{Kaplinghat2016} \\
MW dSph diversity (core-collapse)           & 20--60
   & $\sigma_{\rm eff}/m \sim 20$--$150$ \citep{Correa2021} \\
SDSSJ0946+1006 strong-lens perturber        & $\sim 100$
   & $\sigma_{\rm eff}/m \sim 25$ \citep{Nadler2023} \\
UDG cored rotation curves                   & $\sim 30$--$40$
   & $\sigma_{\rm eff}/m \sim 30$--$100$ \citep{Nadler2023} \\
GD-1 dense substructure                     & $<10$
   & $\sigma_{\rm eff}/m \sim 100$ (needed) \\
\bottomrule
\end{tabularx}

The picture: a mildly-Yukawa cross section running from
$\sim\!100\cmg$ at $\vrel\lesssim 10\kms$ to $\sim\!0.1\cmg$ at
$\vrel\sim 1000\kms$ fits all data.  The corresponding
particle-physics benchmark is a dark force with a $\sim 10$-MeV
mediator; \citet{Kaplinghat2016} gives $m_\chi\sim 15\,{\rm GeV}$,
$m_\phi\sim 17\,{\rm MeV}$; \citet{Correa2021} prefers a sub-GeV DM
scenario with $m_\chi = 0.65\,{\rm GeV}$ and $m_\phi = 0.64\,{\rm
MeV}$.  Both remain consistent with the observational bounds table
above.

The old $\smx\lesssim 0.02\cmg$ bound from cluster shape
\citep{MiraldaEscude2000} has been revisited and superseded ---
lensing observations project through the still-triaxial outer halo,
so the small-projected-radius ellipticity does not directly measure
inner shape and the correct bound is $\smx\lesssim 1\cmg$
\citep{TulinYu2018}.  This is why the modern viable window is much
wider than early SIDM discussions suggested.

%----------------------------------------------------------------------
\section{Connections}
\label{sec:connections}

\subsection{SMBH seeding via gravothermal collapse}

An SIDM halo that undergoes gravothermal core collapse produces a
compact remnant whose mass depends on halo formation redshift and
$\smx$.  The seed's mass at the relativistic instability trigger
(central 3D velocity dispersion $\sim 0.57c$) is a substantial
fraction of the collapsing core mass.

\paragraph{Ultra-strong (uSIDM) sub-component.}
\citet{Pollack2015} propose that a fraction $f\lesssim 0.1$ of the
dark matter has an ``ultra-strong'' cross section $\sigma\gg 1\cmg$
so that $\sigma f\sim 1\cmg$.  The uSIDM component gravothermally
collapses to
$M_{\rm BH} = 0.025\,f\,M_\Delta/[\ln(1+c)-c/(1+c)]$ within
$455.65$ central relaxation times.  For $M_\Delta = 10^{12}\Msun$,
$c=9$, $z_{\rm vir}=15$, the allowed window is
$\sigma f \in [0.34, 0.43]\cmg$, giving a seed of mass sufficient
to grow into the observed $2\times 10^9\Msun$ SMBH of ULAS J1120+0641
at $z=7.085$ under Eddington accretion.

\paragraph{SIDM $+$ baryons.}  \citet{Feng2021} extend this to the
regular SIDM cross section by including a static protogalaxy baryon
potential (fit to Wise+2008 sim).  Baryons shorten the gravothermal
collapse time by factor $\sim 100$, so $\smx\sim 1$--$10\cmg$
(the same value that fits galaxy-scale SIDM observations) produces the
seed for J1205-0000 ($M_{\rm BH} = 2\times 10^9\Msun$, $z=6.7$) in
$\sim 68$ Myr at $z\sim 7.4$.

\paragraph{Intermediate-mass BHs in dSphs.}  \citet{Meshveliani2023}
build a semi-analytic model that predicts IMBH seeds of $0.1$--$1000
\Msun$ in the currently-collapsing MW satellites, following an
$M_{\rm BH}$--$M_{\rm halo}$ relation with the same slope as
extrapolated SMBH-halo scalings but $\sim 2$ dex lower normalisation.
The velocity-dependent ``vd100'' model
($v_{\rm max} = 25\kms$, $\sigma_T^{\rm max} = 60\cmg$) is required to
avoid contradiction with $M_\star\lesssim 3\times 10^4\Msun$ upper
limits on IMBHs in Fornax and Ursa Minor.  Leo I's inferred central
$\sim 3\times 10^6\Msun$ BH remains inconsistent with gravothermal-only
seeding.

\subsection{Atomic dark matter}

See \texttt{dm\_chemistry.tex}.  Atomic DM is a \emph{dissipative}
sector (radiative cooling via dark bremsstrahlung, recombination, and
molecular lines); SIDM is \emph{elastic} (no energy sink other than
gravothermal conduction).  Both soften cores, but by different
mechanisms.  At high densities atomic DM cools and forms compact
objects, while SIDM heats its centre via gravothermal collapse.
Hybrid atomic+SIDM scenarios are physically possible but underexplored.

\subsection{Fuzzy dark matter}

See \texttt{fdm.tex}.  Both FDM and SIDM soften small-scale
structure.  FDM produces a solitonic ground-state core inside every
halo, whose size scales as $r_c\propto 1/(m_\chi^2 M_c)$; SIDM
produces a thermal-equilibrium isothermal core whose radius depends
on local density and $\smx$.  Observational discriminators include
FDM interference-driven granular substructure (absent in SIDM) and
SIDM gravothermal core-collapsed subhaloes (absent in FDM).

%----------------------------------------------------------------------
\section{Summary}
\label{sec:summary}

SIDM is the minimal collisional extension of CDM.  Its
one-parameter-family cross section $\sigma(v)$ modulates a competition
between two regimes: outward heat conduction (core-forming, dominant
at early times and in dilute haloes) and gravothermal catastrophe
(core-collapsing, dominant at late times and in tidally stripped
subhaloes).  Observationally, a velocity-dependent Yukawa cross
section with $\smx\sim 2\cmg$ on dwarf scales and $\lesssim 0.1\cmg$
on cluster scales is preferred over CDM by rotation curves, dwarf
kinematics, and the diversity of low-mass halo properties.  Baryons
substantially compress inner SIDM cores in baryon-dominated galaxies
\citep{Kaplinghat2014,Robles2017}, an effect that is now known to be
essential for producing SMBH seeds \citep{Feng2021}.

Recent numerical work \citep{Gurian2025Core,Schon2025} has focused on
following core collapse into regimes where the fluid approximation
fails, requiring DSMC (KiSS-SIDM) and hybrid particle+fluid (SHH)
methods.  Public simulation compilations \citep{Palubski2025} now
make statistical inference on $\sigma(v)$ from observed substructure
tractable.  The open frontier is joining velocity-dependent SIDM to
baryon feedback, dust-lane hydrodynamics, and the observed
diversity of the Milky Way subhalo population.

%======================================================================
\appendix

\section{Kinetic-Theory Foundation}
\label{app:kinetic}

The gravothermal fluid framework of \S\ref{sec:fluid} sits at the end
of a long derivation chain: Hamiltonian $N$-body $\to$ Liouville
equation $\to$ BBGKY hierarchy $\to$ Boltzmann equation (with closure)
$\to$ fluid (moment) equations $\to$ Navier--Stokes (Chapman--Enskog
closure).  This appendix walks through that chain in enough detail
that the SIDM gravothermal equations
\eqref{eq:balberg_mass}--\eqref{eq:balberg_flux} appear as a specific
application, and the failure mode identified by
\citet{Gurian2025Core} can be located in a specific step.  Presentation
follows the pedagogical order of van den Bosch's ASTR 501 lecture notes
(\emph{Dynamics of Astrophysical Many-Body Systems}).

\subsection{Hamiltonian $N$-body and Liouville's theorem}
\label{app:liouville}

A system of $N$ interacting particles has microstate
$\vec\Gamma = (\vec q_1,\ldots,\vec q_N,\vec p_1,\ldots,\vec p_N)$ in
$6N$-dimensional $\Gamma$-space.  Hamiltonian dynamics evolves
$\vec\Gamma(t)$ via $\dot{\vec q}_i = \partial\mathcal H/\partial\vec
p_i$, $\dot{\vec p}_i = -\partial\mathcal H/\partial\vec q_i$.  The
$N$-body distribution function $f^{(N)}(\vec\Gamma,t)$ describes an
ensemble of microstates.  Probability conservation gives the
\emph{Liouville equation}
%
\begin{equation}
  \frac{\partial f^{(N)}}{\partial t} + \{f^{(N)},\mathcal H\} = 0,
  \quad\text{or equivalently}\quad
  \frac{df^{(N)}}{dt} = 0,
  \label{eq:liouville}
\end{equation}
%
where the \emph{Poisson bracket} of two phase-space functions $A$,
$B$ is
%
\begin{equation}
  \{A,B\} \equiv \sum_{i=1}^{N}\!\left(
      \frac{\partial A}{\partial\vec q_i}\!\cdot\!\frac{\partial B}{\partial\vec p_i}
    - \frac{\partial A}{\partial\vec p_i}\!\cdot\!\frac{\partial B}{\partial\vec q_i}
    \right).
  \label{eq:poisson_bracket}
\end{equation}
%
Substituting Hamilton's equations gives
$\{f^{(N)},\mathcal H\} = \sum_i\bigl(\dot{\vec q}_i\!\cdot\!\partial_{\vec q_i}
+ \dot{\vec p}_i\!\cdot\!\partial_{\vec p_i}\bigr) f^{(N)}$, i.e.\ the
total advective derivative along the Hamiltonian flow --- which is
what makes \eqref{eq:liouville} equivalent to $df^{(N)}/dt = 0$.
That is: the phase-space flow is incompressible, so a comoving
volume element in $\Gamma$-space keeps constant volume.  This is
Liouville's theorem.

\subsection{From Liouville to Boltzmann: the BBGKY hierarchy}
\label{app:bbgky}

The $6N$-dimensional $f^{(N)}$ is unusable in practice.  Define the
reduced $k$-particle distribution function
%
\begin{equation}
  f^{(k)}(\vec w_1,\ldots,\vec w_k, t)
    \equiv \frac{N!}{(N-k)!}\!\int\!\prod_{i=k+1}^{N} d^6\vec w_i\,
           f^{(N)}(\vec w_1,\ldots,\vec w_N,t),
\end{equation}
%
where $\vec w_i \equiv (\vec q_i,\vec p_i)$.  Taking the time
derivative of $f^{(k)}$ and using \eqref{eq:liouville} with the
Hamiltonian split into single-particle + two-body interaction
$U(|\vec q_i - \vec q_j|)$ yields the \emph{BBGKY hierarchy}
(Bogoliubov--Born--Green--Kirkwood--Yvon):
%
\begin{equation}
  \frac{\partial f^{(k)}}{\partial t} = \{\mathcal H^{(k)}, f^{(k)}\}
   + \sum_{i=1}^{k}\!\int\! d^3\vec q_{k+1}\,d^3\vec p_{k+1}\;
     \frac{\partial U(|\vec q_i - \vec q_{k+1}|)}{\partial \vec q_i}
     \cdot\frac{\partial f^{(k+1)}}{\partial\vec p_i}.
  \label{eq:bbgky}
\end{equation}
%
The evolution of $f^{(k)}$ depends on $f^{(k+1)}$: an unclosed
hierarchy of $N$ coupled equations.  Nothing has been assumed except
Hamiltonian dynamics.

\paragraph{Mayer cluster expansion.}  Introduce correlation functions
via
%
\begin{align}
  f^{(2)}(1,2) &= f(1)f(2) + g(1,2), \\
  f^{(3)}(1,2,3) &= f(1)f(2)f(3) + f(1)g(2,3) + \text{perm.} + h(1,2,3).
\end{align}
%
$g(1,2)$ measures two-body correlations; $h$ measures three-body
irreducible correlations.  Uncorrelated particles have
$g = h = 0$.

\paragraph{Closure: molecular chaos (Stosszahlansatz).}  For a
\emph{dilute gas with short-range} two-body interactions (potential
$U$ has support only in a small collision volume $r_{\rm coll}$), one
closes \eqref{eq:bbgky} by assuming velocities entering a collision are
uncorrelated, $f^{(2)}(\vec x,\vec x,\vec p_1,\vec p_2) = f^{(1)}(\vec
x,\vec p_1)\,f^{(1)}(\vec x,\vec p_2)$ (identity of positions since
the interaction is localized).  The $k=1$ equation of \eqref{eq:bbgky}
then closes on itself as the \emph{Boltzmann equation}
%
\begin{equation}
  \frac{df}{dt}
   \equiv \frac{\partial f}{\partial t}
        + \vec v\cdot\nabla_{\!\vec x} f
        - \nabla\Phi\cdot\nabla_{\!\vec v} f
   = I[f],
  \label{eq:boltzmann}
\end{equation}
%
with (elastic, two-body) collision integral
%
\begin{equation}
  I[f] = \!\int\! d^3\vec p_2\,d^3\vec p_1'\,d^3\vec p_2'\;
    \omega(\vec p_1',\vec p_2'|\vec p_1,\vec p_2)
    \bigl[f(\vec p_1')f(\vec p_2') - f(\vec p_1)f(\vec p_2)\bigr].
  \label{eq:collision_integral}
\end{equation}
%
Here $\omega$ is the differential cross section, $f\equiv f^{(1)}$
throughout, and the two terms describe replenishing (gain) and
depleting (loss) collisions.  Molecular chaos secretly introduces an
arrow of time --- the Loschmidt paradox --- but works empirically.

\paragraph{Closure: uncorrelated momenta (Vlasov--Poisson).}  For
gravitational $N$-body, interactions are long-range and no
$r_{\rm coll}$ exists.  The closure that works instead is stronger
than molecular chaos: assume momenta at all pairs are uncorrelated,
$g(1,2) = 0$ everywhere.  The two-body interaction integral in
\eqref{eq:bbgky} then combines with the streaming term to give the
mean-field potential $\Phi(\vec x)$ sourced by $\rho = m\!\int f\,d^3\vec
p$ via Poisson.  Equation \eqref{eq:boltzmann} reduces to the
\emph{collisionless Boltzmann equation} (CBE, or Vlasov--Poisson):
%
\begin{equation}
  \frac{\partial f}{\partial t}
    + \vec v\cdot\nabla_{\!\vec x} f
    - \nabla\Phi\cdot\nabla_{\!\vec v} f = 0,
  \qquad \nabla^2\Phi = 4\pi G\!\int m f\,d^3\vec v.
  \label{eq:cbe}
\end{equation}
%
This governs pure CDM $N$-body simulations.  In a collisionless
system, $g(1,2) = 0$ is preserved exactly.  \emph{SIDM sits between
these two closures}: elastic short-range scattering \emph{plus}
gravitational Poisson.  The Boltzmann equation
\eqref{eq:boltzmann}--\eqref{eq:collision_integral} with $\Phi$
determined by Poisson is the master equation for SIDM in
phase-space.

\subsection{Moment equations of the Boltzmann equation}
\label{app:moments}

Take a scalar function $Q(\vec v)$ and integrate the Boltzmann
equation multiplied by $Q$ over velocity space.  Define
%
\begin{equation}
  \langle Q\rangle(\vec x, t) \equiv \frac{1}{n(\vec x,t)}\!\int Q(\vec v) f\,d^3\vec v,
  \qquad
  n = \!\int f\,d^3\vec v.
\end{equation}
%
For $Q$ a \emph{collisional invariant} --- conserved by every
collision, i.e.\ $Q(\vec p_1) + Q(\vec p_2) = Q(\vec p_1') + Q(\vec
p_2')$ --- one can show $\int Q\,I[f]\,d^3\vec v = 0$.  The three
collisional invariants for elastic scattering are
$Q\in\{1,m\vec v,\tfrac{1}{2}mv^2\}$: particle number, momentum,
energy.  Their moment equations coincide with the collisionless
moment equations \emph{despite} the microscopic collisions being
present --- collisionality only shows up in higher moments and in the
stress tensor.

Standard manipulations (integration by parts in $\vec v$-space with
$f\to 0$ at infinity, and the divergence trick for the
$\vec v\cdot\nabla f$ streaming term) yield the master moment equation
%
\begin{equation}
  \frac{\partial}{\partial t}(n\langle Q\rangle)
   + \frac{\partial}{\partial x_i}
     \bigl[n\langle Q\,v_i\rangle\bigr]
   + \frac{\partial\Phi}{\partial x_i}
     n\Bigl\langle\frac{\partial Q}{\partial v_i}\Bigr\rangle = 0.
  \label{eq:master_moment}
\end{equation}

\paragraph{$Q=m$: continuity.}
$\partial_t\rho + \partial_i(\rho u_i) = 0$
with $\rho = mn$ and macroscopic velocity $u_i = \langle v_i\rangle$.

\paragraph{$Q=mv_j$: momentum.}  Split microscopic velocity
$\vec v = \vec u + \vec w$ into macroscopic streaming $\vec u$ and
random peculiar $\vec w$ (so $\langle\vec w\rangle = 0$).  Then
$\rho\langle v_i v_j\rangle = \rho u_i u_j + \rho\langle w_i w_j\rangle$.
Introducing the \emph{stress tensor}
%
\begin{equation}
  \sigma_{ij} \equiv -\rho\langle w_i w_j\rangle,
  \label{eq:stress_general}
\end{equation}
%
\eqref{eq:master_moment} for $Q=mv_j$ reads
%
\begin{equation}
  \frac{\partial u_j}{\partial t} + u_i\partial_i u_j
    = \frac{1}{\rho}\partial_i\sigma_{ij} - \partial_j\Phi.
  \label{eq:momentum_equation}
\end{equation}
%
The stress tensor is where microscopic physics lives.

\paragraph{$Q=\tfrac{1}{2}mv^2$: energy.}  Define specific internal
energy $\varepsilon = \tfrac{1}{2}\langle w^2\rangle$.  The energy
equation reads
%
\begin{equation}
  \rho\frac{d\varepsilon}{dt}
   = -P\,\partial_k u_k + \mathcal V - \partial_k F_{{\rm cond},k},
  \label{eq:energy_equation}
\end{equation}
%
with viscous dissipation rate $\mathcal V$ and conductive heat flux
$\vec F_{\rm cond}$.

\subsection{The Chapman--Enskog closure}
\label{app:CE}

Equations \eqref{eq:momentum_equation}--\eqref{eq:energy_equation}
are unclosed: $\sigma_{ij}$, $\mathcal V$, and $F_{\rm cond}$ all
depend on higher moments of $f$.  The Chapman--Enskog procedure is a
perturbation expansion in the \emph{Knudsen number}
%
\begin{equation}
  \mathrm{Kn} \equiv \frac{\lambda_{\rm mfp}}{L_{\rm macro}},
\end{equation}
%
where $L_{\rm macro}$ is the shortest gradient length in the fluid
variables.  Small Kn = fluid regime; large Kn = free-streaming
(collisionless) regime.

\paragraph{Zeroth order.}  Locally set $I[f] = 0$, forcing $f =
f^{\rm MB}$ a local Maxwell--Boltzmann distribution with parameters
$\rho(\vec x,t)$, $\vec u(\vec x,t)$, $T(\vec x,t)$:
%
\begin{equation}
  f^{\rm MB}(\vec v)
   = \frac{n}{(2\pi m k_B T)^{3/2}}\,
     \exp\!\left[-\frac{|\vec v-\vec u|^2\,m}{2k_B T}\right].
  \label{eq:maxwellian}
\end{equation}
%
At this order $\sigma_{ij}^{(0)} = -P\,\delta_{ij}$ with
$P = n k_B T$ (ideal gas law): \emph{Euler} equations, no viscosity,
no conduction.

\paragraph{First order.}  Expand $f = f^{\rm MB}\,(1 + \phi_1 + \cdots)$
with $\phi_1 \propto \mathrm{Kn}$.  Substituting into the Boltzmann
equation and linearising the collision integral around
$f^{\rm MB}$ gives a linear integral equation for $\phi_1$ whose
solution reveals
%
\begin{equation}
  \sigma_{ij} = -P\,\delta_{ij}
    + \mu\!\left(\partial_i u_j + \partial_j u_i
                - \tfrac{2}{3}\delta_{ij}\partial_k u_k\right)
    + \eta\,\delta_{ij}\partial_k u_k,
  \label{eq:navier_stokes_stress}
\end{equation}
%
%
\begin{equation}
  \vec F_{\rm cond} = -\kappa\,\nabla T.
  \label{eq:fourier_law}
\end{equation}
%
The three transport coefficients $\mu$ (shear viscosity),
$\eta$ (bulk viscosity), and $\kappa$ (thermal conductivity) are
computed as integrals of collision kernels over Maxwellian velocity
distributions.  For elastic \emph{hard-sphere} scattering with total
cross section $\sigma$ and particle mass $m$
(Chapman--Enskog result, see any kinetic-theory textbook):
%
\begin{align}
  \mu &= \frac{5}{16}\,\frac{1}{\sigma}\sqrt{\frac{m k_B T}{\pi}},
    &&\text{(shear viscosity)},\\
  \kappa &= \frac{75\,k_B}{64\,\sigma}
             \sqrt{\frac{k_B T}{\pi m}},
    &&\text{(thermal conductivity)}.
  \label{eq:CE_transport}
\end{align}
%
Neither depends on $n$: at higher density, more particles carry the
flux but each carries less because $\lambda_{\rm mfp}$ shrinks.  Bulk
viscosity $\eta = 0$ for a monatomic gas.

Substituting \eqref{eq:navier_stokes_stress}--\eqref{eq:fourier_law}
into \eqref{eq:momentum_equation} and \eqref{eq:energy_equation}
yields the \emph{Navier--Stokes equations}: the standard collisional
fluid model.  Zero transport coefficients recovers Euler.

\paragraph{Higher orders (Burnett, super-Burnett).}  Expanding to
$\mathrm{Kn}^2$ adds gradient-of-gradient terms with typically
non-physical (backwards-in-time) solutions.  In practice one stops at
Navier--Stokes.

\subsection{Thermal conduction: from mean free path to Chapman--Enskog}
\label{app:conduction}

Thermal conduction is the heat-flux channel that drives the
gravothermal catastrophe (\S\ref{sec:regimes}), so it deserves a
dedicated derivation.  Two approaches, matched at the end.

\subsubsection{Elementary mean-free-path estimate}

Consider a gas with a small temperature gradient $dT/dz$ along the
$z$-axis, no bulk flow.  A particle at height $z$ typically had its
last collision at height $z - \lambda_{\rm mfp}\cos\theta_v$, where
$\theta_v$ is the angle between its velocity and $\hat z$.  It carries
the average energy of that location:
$\bar\varepsilon = \tfrac{3}{2}k_B T(z - \lambda_{\rm mfp}\cos\theta_v)$.
The net upward heat flux is the difference between upward-moving and
downward-moving particles integrated over the Maxwellian:
%
\begin{equation}
  F_z \;=\; -n\,\langle v_z\bigl(\bar\varepsilon(z-\lambda_{\rm mfp}\cos\theta_v)
        -\bar\varepsilon(z)\bigr)\rangle
    \;\simeq\; -n\,\langle v_z^2\rangle\,\lambda_{\rm mfp}\,
        \frac{3k_B}{2}\,\frac{dT}{dz}.
\end{equation}
%
Using $\langle v_z^2\rangle = k_B T/m$ and
$\lambda_{\rm mfp} = 1/(n\sigma)$ gives
%
\begin{equation}
  F_z = -\kappa_{\rm mfp}\frac{dT}{dz},
  \qquad
  \kappa_{\rm mfp}
    = \tfrac{3}{2}\,\frac{k_B}{\sigma}\sqrt{\frac{k_B T}{m}}
    \sim n\,c_v\,\bar v\,\lambda_{\rm mfp}.
  \label{eq:kappa_mfp}
\end{equation}
%
Compact form: $\kappa \sim n c_v \bar v\lambda_{\rm mfp}$, with $c_v$
the specific heat per particle, $\bar v$ the thermal speed,
$\lambda_{\rm mfp}$ the mean free path.  The factor of $n$ in
$n\lambda_{\rm mfp} = 1/\sigma$ cancels: increase $n$ and each
particle carries less flux per collision, exactly compensating.
Consequence: $\kappa$ is independent of density.

\subsubsection{Full Chapman--Enskog calculation}

The mean-free-path estimate misses order-unity numerical constants
that come from the Maxwellian velocity distribution and from the
velocity dependence of the mean free path itself (fast particles have
shorter $\lambda_{\rm mfp}$).  The full Chapman--Enskog result
requires solving a linear integral equation for the first-order
correction $\phi_1$ to the local Maxwellian.

\paragraph{Setup.}  Expand
$f = f^{\rm MB}\bigl(1 + \phi_1(\vec c) + \cdots\bigr)$
with $\vec c \equiv \vec v - \vec u$ the peculiar velocity.  Choose
the ansatz relevant for thermal conduction (no bulk flow gradient):
%
\begin{equation}
  \phi_1(\vec c) = -\vec A(\vec c)\cdot\nabla\ln T,
\end{equation}
%
so $\phi_1$ is linear in $\nabla T$ and vanishes in equilibrium.
Substitute into the linearised Boltzmann equation.  The streaming
side gives, after some algebra,
%
\begin{equation}
  f^{\rm MB}\!\left(\frac{mc^2}{2 k_B T} - \tfrac{5}{2}\right)\vec c\cdot\nabla\ln T,
  \label{eq:CE_source}
\end{equation}
%
which must be balanced by the linearised collision integral acting on
$\vec A$.

\paragraph{Sonine expansion.}  The angular structure of \eqref{eq:CE_source}
forces $\vec A = A(c^2)\,\vec c$ for some scalar $A$.  Expand $A(c^2)$
in Sonine (associated Laguerre) polynomials:
%
\begin{equation}
  A(c^2) = \sum_{p=0}^\infty a_p\,S_{3/2}^{(p)}\!\!\left(\frac{mc^2}{2k_B T}\right)\!,
\end{equation}
%
where $S_{3/2}^{(p)}$ is the Sonine polynomial that diagonalises the
linearised collision operator for a Maxwellian equilibrium.
Truncating at first Sonine ($p=0$) gives the leading result;
subsequent orders converge rapidly.

\paragraph{Result.}  Solving for $a_0$ using the orthogonality of
Sonine polynomials against the collision kernel gives, for elastic
hard-sphere scattering with total cross section $\sigma$
(Chapman \& Cowling 1970, \S 10.21):
%
\begin{equation}
  \kappa = \frac{75\,k_B}{64\,\sigma}\sqrt{\frac{k_B T}{\pi m}}.
  \label{eq:kappa_CE_full}
\end{equation}
%
Comparing to \eqref{eq:kappa_mfp}:
$\kappa_{\rm CE}/\kappa_{\rm mfp} = 25\pi/64 \approx 1.23$.  The
mean-free-path estimate is off by only $\sim 25\%$: the Sonine
expansion is efficient because the Maxwellian distribution is close
to a $\delta$-function in $c^2$ for the moments that matter.

\subsubsection{Higher-Sonine corrections}

Adding the $p=1$ Sonine term corrects \eqref{eq:kappa_CE_full} by
about $1\%$; $p=2$ by less than $10^{-3}$.  Rigorous bounds
(Uhlenbeck) show the first-Sonine result is a strict lower bound and
that the true $\kappa$ lies within $2\%$ for hard spheres.  The
canonical value \eqref{eq:kappa_CE_full} is used throughout the SIDM
literature.

\subsubsection{From Chapman--Enskog to gravothermal SIDM}

For an SIDM halo we replace $k_B T \to m\nu^2$ (SIDM is a
self-gravitating gas whose ``temperature'' is the 1D velocity
dispersion in kinetic units).  Substituting into
\eqref{eq:kappa_CE_full} and expressing $\sigma$ per unit mass:
%
\begin{equation}
  \kappa_{\rm SIDM,SMFP}
    = \frac{75}{64\sqrt\pi}\,\frac{m\nu}{\sigma}
    \simeq 2.1\,\frac{m\nu k_B}{\sigma}\,\frac{1}{k_B}
    = 2.1\,\frac{\nu\,k_B}{\smx}.
  \label{eq:kappa_SIDM_SMFP}
\end{equation}
%
This matches Eq.~\eqref{eq:balberg_flux}'s SMFP limit within the
$b = 25\sqrt\pi/32$ conductivity factor of \citet{Balberg2002} up to
$O(1)$ manipulations.  The corresponding heat flux from Fourier's law
is
%
\begin{equation}
  \vec F_{\rm cond}^{\rm SMFP}
   = -\kappa_{\rm SIDM,SMFP}\,\nabla\!\bigl(m\nu^2/k_B\bigr)
   \propto -\frac{\nu}{\sigma}\,\nabla\nu^2,
\end{equation}
%
which is exactly what appears in the fluid-model luminosity
\eqref{eq:balberg_flux}.

\subsubsection{The LMFP heat flux: gravothermal conduction}

Chapman--Enskog assumes $\lambda_{\rm mfp}\ll H$.  In the outer halo
(and initially, before any core forms), $\lambda_{\rm mfp}\gg H$ and
the fluid picture manifestly fails: a particle traverses many scale
heights before scattering, so it samples the temperature far from its
starting location.  The correct LMFP heat flux is derived in
\citet{LyndenBell1980} and \citet{Balberg2002}: particles travel a
gravitational scale height
$H \equiv (\nu^2/4\pi G\rho)^{1/2}$ between scatterings, and each
scattering communicates the temperature difference over that scale.
Effective ``mean free path for heat'' becomes $H$, and the effective
conductivity is
%
\begin{equation}
  \kappa_{\rm SIDM,LMFP}
    \simeq 0.27\,C\,\frac{\rho\,\nu^3\,\sigma\,k_B}{Gm^2}
    \propto \rho\,\sigma
    \qquad (\lambda_{\rm mfp}\gg H).
  \label{eq:kappa_SIDM_LMFP}
\end{equation}
%
Note the opposite dependence on $\sigma$: SMFP conductivity is
$\propto 1/\sigma$ (more collisions $\to$ slower diffusion), LMFP
conductivity is $\propto\sigma$ (more collisions $\to$ more energy
transferred between distant zones).  This is why the interpolated
form
$\kappa^{-1} = \kappa_{\rm SMFP}^{-1} + \kappa_{\rm LMFP}^{-1}$
(used by \citealt{Balberg2002,Pollack2015,Meshveliani2023}) captures
both regimes: in SMFP the term $\kappa_{\rm SMFP}$ dominates, in LMFP
$\kappa_{\rm LMFP}$ does.  The interpolation is not a Chapman--Enskog
result --- it is an ansatz calibrated against $N$-body simulations,
with the $C\simeq 0.75$ (or $0.42$--$0.84$ depending on source)
fudge factor absorbing the mismatch.

\subsubsection{The IMFP problem}

The $\mathrm{Kn}\sim 1$ region is where the interpolation breaks
down.  Chapman--Enskog says $\kappa$ diverges as
$\kappa_{\rm SMFP}\propto 1/\sigma$; LMFP says $\kappa_{\rm LMFP}
\propto\rho\sigma$; these two functional forms only match in a
neighbourhood of $\rho\sigma H = O(1)$.  \citet{Gurian2025Core} DSMC
runs show that in the IMFP shell:
%
\begin{itemize}
  \item The velocity distribution is anisotropic and
        non-Maxwellian in its radial component: the radial
        distribution has a mode near zero and an extended,
        asymmetric outward tail --- boil-off particles that have not
        thermalized --- while the angular components stay closer to
        Maxwellian (Fig.~3 of \citealt{Gurian2025Core}; the
        radial-dispersion enhancement is visible in their Fig.~1).
  \item Ordinary Fourier's law $F = -\kappa\nabla T$ gives the wrong
        \emph{sign} in some sub-shells (energy against the gradient),
        because non-local energy transport dominates.
  \item The temperature-weighted entropy-production rate
        $\dot Q = \int dM\,\tfrac13\langle v^2\rangle\,(ds_1/dt)_M$
        (which pure conduction forces to zero when integrated over
        the halo --- the integrand telescopes to boundary
        luminosities) settles at a positive asymptote: entropy is
        being \emph{produced}, not transported, which conduction
        cannot do and viscous dissipation can.
\end{itemize}
%
Fixing this within a fluid framework requires modelling $\sigma_{ij}$
with a genuine viscous term \eqref{eq:navier_stokes_stress} (currently
neglected in gravothermal models), computing higher-Sonine
corrections, or matching to a solution of the actual Boltzmann
equation via DSMC (Bird 1994).

\subsection{Application: gravothermal fluid model for SIDM}
\label{app:sidm_fluid}

The \citet{Balberg2002} gravothermal fluid model of SIDM combines
three inputs:
%
\begin{enumerate}
  \item Momentum and mass conservation from the moment equations
        \eqref{eq:momentum_equation}, adapted to spherical symmetry
        with hydrostatic equilibrium
        $\partial_r(\rho\nu^2) = -GM\rho/r^2$.
  \item Chapman--Enskog thermal conductivity from
        \eqref{eq:CE_transport}, rewritten in terms of the SIDM
        cross section $\smx$ and 1D velocity dispersion $\nu$
        (using $k_B T \to m\nu^2$):
        %
        \begin{equation}
          \kappa_{\rm SMFP}
            = \frac{75}{64}\,\frac{\nu\,k_B}{\sigma}
            \simeq 2.1\,\frac{\nu k_B}{\sigma},
          \label{eq:kappa_smfp}
        \end{equation}
        %
        matching Eq.~\eqref{eq:balberg_flux}'s SMFP limit up to the
        factor $75/64\simeq 1.17$.
  \item An ad hoc interpolation to the LMFP regime where
        Chapman--Enskog fails ($\mathrm{Kn}\gg 1$), giving the
        conducting-flux ansatz Eq.~\eqref{eq:balberg_flux}.
\end{enumerate}
%
The energy equation \eqref{eq:energy_equation} closes with
$\vec F_{\rm cond}$ from \eqref{eq:fourier_law} and Fourier's law.
Neglecting viscous dissipation ($\mathcal V = 0$) --- which is the
standard fluid ansatz for the gravothermal problem --- yields
\eqref{eq:balberg_energy}.

The full derivation of $\kappa_{\rm SMFP}$ for gravitating hard
spheres, including the numerical constant $b = 25\sqrt\pi/32\simeq
1.0$ that appears in \citet{Balberg2002}, is in the classic kinetic
theory literature (Chapman \& Cowling 1970; Huang 1987).

\subsection{Where the fluid model fails}
\label{app:CE_breakdown}

Chapman--Enskog requires $\mathrm{Kn}\lesssim 1$ everywhere.  In an
SIDM halo, $\mathrm{Kn}$ varies by many orders of magnitude between
the outer halo ($\lambda_{\rm mfp}\gg H$, free-streaming) and the
collapsing centre ($\lambda_{\rm mfp}\ll H$, fluid-like).  At any
given time there is a shell where $\mathrm{Kn}\sim 1$ --- the
intermediate-mean-free-path (IMFP) region.

\citet{Gurian2025Core} used direct simulation Monte Carlo
(KiSS-SIDM) to show that in the IMFP region:
%
\begin{itemize}
  \item The velocity distribution is \emph{anisotropic}
        ($\langle w_r^2\rangle \neq \langle w_\perp^2\rangle$), not
        locally Maxwellian.
  \item Non-Maxwellian tails in the radial component develop
        (``boil-off'' particles).
  \item The entropy-production diagnostic $\hat Q$ turns positive,
        implying the fluid-model heat flow is going against the
        temperature gradient --- a direct signature of Chapman--Enskog
        breakdown.
\end{itemize}
%
Physically, once $\mathrm{Kn}\sim 1$ the mean free path becomes
comparable to the shell thickness, so particles no longer thermalise
locally.  The Chapman--Enskog expansion is not just quantitatively
inaccurate; its zeroth-order local-Maxwellian assumption is
qualitatively wrong.  Fixing this requires either (i) computing
higher-order terms in the Chapman--Enskog expansion (rarely
practical), (ii) direct simulation Monte Carlo (KiSS-SIDM), or
(iii) hybrid particle+fluid schemes that transition between regimes
smoothly (SHH; \citealt{Schon2025}).

\bibliography{references}

\end{document}
